\begin{lemma}[The germ ring is a domain; zero sets are thin]
\label{lag:germ-domain}
\lean{TranscendentalNakamaye.Pending.RiemannExtension.dense_compl_zeroSet,TranscendentalNakamaye.Pending.RiemannExtension.isDomain_holGerm}
\leanok
If \(h\) is holomorphic on a connected open set \(U\subset\mathbb C^d\) and does
not vanish identically near some point of \(U\), then \(h^{-1}(0)\) has empty
interior in \(U\). The ring \(\mathcal O_d\) is an integral domain.
\end{lemma}
\begin{proof}
\leanok
The set of points near which \(h\) vanishes identically is open, and it is closed
in \(U\) by the identity theorem (Taylor coefficients at a limit point vanish).
It is not all of \(U\), so by connectedness it is empty; thus \(h^{-1}(0)\)
contains no open set. If \(fg=0\) near \(0\) with \(f,g\) holomorphic on a ball
\(B\) and \(f\) not identically zero near \(0\), then \(g=0\) on the dense set
\(B\setminus f^{-1}(0)\), hence on \(B\) by continuity.
\end{proof}
